% Summary of the Experiments and Evaluation chapter
\section{Summary}
\label{sec:exp-summary}

% PERFORMANCE

The data show that objective performance is higher in \ac{cv}, while improvement rates are higher in \ac{pv}. This is explained mainly by the fact that \ac{cv} has a higher baseline skill in \ac{uml}, which also reduces the potential for improvement among its participants. 

The results show consistently higher scores across activities in \ac{cv}. This may be related not only to higher starting point, but also to behavior patterns. Students in \ac{cv} engaged with and retried quizzes more often than students in \ac{pv}.

On the other hand, students in \ac{pv} engaged more with the theoretical content. This indicates differences in interaction patterns between the groups - an exploratory approach in \ac{pv} and a goal-oriented approach in \ac{cv}.

The textual responses and questionnaire results indicate better overall reception and higher enthusiasm than \ac{cv} across most subjective metrics. However, \ac{pv} also shows a decrease in interest in software modeling after the end of the course, suggesting a potential risk of shifting motivation away from intrinsic interest if not designed carefully.

Responses in \ac{cv} are more variable, indicating a polarizing effect of the design. \ac{pv} responses are more consistent, suggesting greater accessibility.

Overall, competitive design appears more suitable for already skilled and interested students, while playful design is more accessible to a broader audience.

Students expressed general interest in e-learning materials and their use at \ac{fiit}. The reception of this course suggests further experiments on the topic of e-learning at \ac{fiit} are justified.

% \ac{cv} contains a group of students who had high scores in the intro quiz, but seemed to lose interest and ended up with significantly lower scores in the final quiz. This further demonstrates the observed division in \ac{cv}, where 

% We can again observe a division in \ac{cv}, where a number of students seemed to lose interest during the course, which did not happen in \ac{pv}.

% ENGAGEMENT

% The textual responses in \ac{pv} indicate better overall reception and higher enthusiasm than \ac{cv}, which is supported by higher satisfaction and engagement scores across most metrics.

% However, \ac{pv} shows a decrease in interest in software modeling, suggesting a potential risk of shifting motivation away from intrinsic interest if not designed carefully.

% Responses in \ac{cv} are more variable, indicating a polarizing effect, whereas \ac{pv} responses are more consistent, suggesting greater accessibility.

% Overall, competitive design appears more suitable for already skilled and interested students, while playful design is more accessible to a broader audience.

% FIIT PROSPECTS

% Students expressed general interest in e-learning materials and their use at \ac{fiit}. The reception of this course suggests further experiments on the topic of e-learning at \ac{fiit} should be conducted with a preference for playfulness over competition. 


\textbf{RQ1:} Does competitive course design lead to higher learning performance compared to playful design?

\textbf{Answer to RQ1:} Participants in \ac{cv} achieved consistently higher scores across exercises, categories and assessments. However, participants in \ac{pv} achieved higher improvement rates, which can be largely attributed to lower baseline performance and more room for improvement. Although results suggest a slight performance advantage in \ac{cv}, both approaches lead to comparable learning outcomes overall.

\textbf{RQ2:} How do playful and competitive course designs influence students’ behavioral engagement?

\textbf{Answer to RQ2:} The designs led to different behavioral patterns. Participants in \ac{pv} interacted more with the theoretical content, while participants in \ac{cv} had more quiz interactions and attempts. This suggests a tendency towards more exploratory behavior in \ac{pv}, and a more goal-oriented attitude in \ac{cv}.

\textbf{RQ3:} How do students evaluate playful and competitive course designs in terms of emotional engagement, satisfaction, and motivation?

\textbf{Answer to RQ3:} Participants in \ac{pv} expressed higher enthusiasm, satisfaction and engagement in nearly all metrics as well as textual responses.

However, they also showed a drop in interest in software modeling itself between the introduction and feedback questionnaires. This may indicate a shift in motivation or a difference in how the topic is perceived when presented in a playful manner, although the effect is not statistically significant.

Responses in \ac{cv} exhibited higher variability, indicating that competition may have a polarizing effect on the students' perceptions. \ac{pv} had more consistent responses and proved to be more accessible in general. 
% A number of students in the competitive design 

\textbf{RQ4:} What are students’ perceptions of the usefulness and future potential of e-learning materials at \ac{fiit}?

\textbf{Answer to RQ4:} Students perceived the course as a valuable addition to the software modeling subject and expressed interest in encountering more e-learning materials during their studies at \ac{fiit}. Higher general interest was expressed regarding playful materials than competitive ones.
